%%
%% fall-lecture17.tex
%% 
%% Made by Alex Nelson <pqnelson@gmail.com>
%% Login   <alex@lisp>
%% 
%% Started on  2026-11-05T07:44:55-0800
%% Last update 2026-11-05T07:44:55-0800
%% 

\lecture{}

We will look at a couple more definitions in Thermodynamics, then we
will start studying ideal gases (and perfect gases).

\subsection{Higher-order derivatives}

\begin{node}
So far, we have been talking about Thermodynamic quantities which are
first-order derivatives. But there are quantities which are
second-order derivatives.
\end{node}

\begin{definition}
The \define{(Dimensionless) Fundamental Derivative of Gas Dynamics}
(or \define{Landau Derivative}) is defined to be
\begin{equation*}
\Gamma:=\frac{c^{4}}{2V^{3}}\left(\frac{\partial^{2}V}{\partial p^{2}}\right)_{S},
\end{equation*}
which tells us the curvature of isentropes in $p$-$V$ space.
\end{definition}

\begin{node}
The curvature of $V$ with respect to $p$ is precisely $\Gamma$. But,
we can compute the curvature of $p$ with respect to $V$ as:
\begin{subequations}
  \begin{align}
\left(\frac{\partial^{2}p}{\partial V^{2}}\right)_{S}
&= \frac{\partial}{\partial V}\left(\frac{\partial p}{\partial V}\right)_{S}\\
&= \frac{\partial}{\partial V}\left[\frac{1}{(\partial V/\partial p)_{S}}\right]\\
&= -\left[\frac{1}{(\partial V/\partial p)_{S}}\right]^{2}\left(\frac{\partial^{2}V}{\partial p^{2}}\right)_{S}\\
&= \frac{-c^{4}}{V^{4}}\left(\frac{\partial^{2}V}{\partial p^{2}}\right)_{S}\quad\mbox{by (\S\ref{eq:lec16:fall:partial-V-partial-p})}\\
&= \frac{V^{3}}{2c^{2}}\left(\frac{\partial^{2}V}{\partial p^{2}}\right)_{S}
  \end{align}
\end{subequations}
So we can describe the curvature of $p$ with respect to $V$ using
$\Gamma$. For other forms, see Thompson~\cite{thompson1971fundamental}.
\end{node}

\begin{node}
Why are these useful? How do they come about? Suppose we have some
Thermodynamic system and it starts at $(V_{1},p_{1})$:
\begin{center}
\includegraphics{img/img.30}
\end{center}
Then we move along the isentrope with constant entropy $S=S_{1}$ to
the point $(V_{2},p_{2})$:
\begin{center}
\includegraphics{img/img.31}
\end{center}
Then we move to point $(V_{3},p_{2})$ via an isobaric process [holding
pressure constant] to the isentrope with entropy
$S_{2}$ by increasing the volume:
\begin{center}
\includegraphics{img/img.32}
\end{center}
We have
\begin{equation}
\Delta V = V_{3}-V_{1},\quad\mbox{and}\quad\Delta p =p_{3}-p_{1}=p_{2}-p_{1}
\end{equation}
and
\begin{equation}
\Delta S = S_{2}-S_{1}.
\end{equation}
Looking at the ratio $(\Delta V)/V$, we can Taylor expand it
\begin{multline}
\frac{\Delta V}{V} = \left[\frac{1}{V}\left(\frac{\partial V}{\partial p}\right)_{S}\Delta p
+ \frac{1}{2!}\frac{1}{V}\left(\frac{\partial^{2} V}{\partial p^{2}}\right)_{S}(\Delta p)^{2}
+\bigOh(\Delta p^{3})\right] + \\
+\left[\frac{1}{V}\left(\frac{\partial V}{\partial S}\right)_{p}\Delta S + \bigOh(\Delta S^{2})\right]
\end{multline}
where the first line describes the isentropic process, and the second
line describes the isobaric process. We can plug in some of our
Thermodynamic definitions to find
\begin{subequations}
  \begin{align}
\frac{\Delta V}{V}
&= -\kappa_{S}\Delta p + c\Gamma\frac{(\Delta p)^{2}}{2} + \bigOh(\Delta p^{3})+\frac{1}{V}\frac{T\alpha V}{c_{p}}\Delta S + \bigOh(\Delta S^{2})\\
&=\underbrace{\frac{-\Delta p}{\rho c^{2}} + \Gamma\left(\frac{\Delta p}{\rho c}\right)^{2}}_{\text{isentropic compression}}
    + \underbrace{G\frac{T\,\Delta S}{c^{2}}}_{\text{external heating}}+ \bigOh(\Delta p^{3}) + \bigOh(\Delta S^{2})
  \end{align}
\end{subequations}
where the first two terms describe isentropic compression, the third
term describes external heating from non-isentropic compression. Here
we introduced
\begin{equation}
G := V\cdot\left(\frac{\partial p}{\partial\rho}\right)_{V}
\end{equation}
which is called the \define{Nondimensional Thermal Expansion Coefficient}.
For gases usually $0<G<2/3$ and for liquids $G\sim2$.
\end{node}

\subsection{Ideal gas Thermodynamics}

\begin{assume}
We will have three assumptions:
\begin{enumerate}
\item Interparticle interactions are negligible
\item We have a dilute gas
\item The gas is at equilibrium.
\end{enumerate}
\end{assume}

\begin{node}
We can describe the idea gas Law using statistical mechanics. Recall
the Maxwell distribution, where $\vec{u}$ is the mean velocity:
\begin{equation}
f_{0}(\vec{c})=\left(\frac{m}{2\pi\Boltzmann}\right)\exp\left(\frac{-m}{2RT}\|\vec{c}-\vec{u}\|^{2}\right).
\end{equation}
We can always change reference frames to make $\vec{u}=0$ (thanks to
Galilean relativity!).
\end{node}

\begin{node}
Temperature is defined as the kinetic energy of these molecules
\begin{equation*}
\iiint_{\RR^{3}}\frac{1}{2}m(c_{1}^{2}+c_{2}^{2}+c_{3}^{2})f_{0}(\vec{c})\,\D\vec{c}=3\left(\frac{1}{2}RT\right).
\end{equation*}
\end{node}

\begin{node}
Pressure is just the flux of momentum per unit area:
\begin{subequations}
  \begin{align}
p &=\frac{1}{3}\rho\iiint_{\RR^{3}}(c_{1}^{2}+c_{2}^{2}+c_{3}^{2})f_{0}(\vec{c})\,\D\vec{c}\\
&=\frac{1}{3}\rho(3RT)\quad\mbox{by previous calculation}\\
&=\rho RT
  \end{align}
\end{subequations}
where $R$ is the [nondimensionalized] universal gas constant,
\begin{equation}
R = \frac{\widetilde{R}}{m}
\end{equation}
where $m$ is the molar mass, and $\widetilde{R}\approx8.314~\mathrm{J}\cdot\mathrm{K}^{-1}\cdot\mathrm{mol}^{-1}$.
\end{node}

\begin{node}
Observe from the definition of pressure $p=\rho RT$, we can rewrite
this as
\begin{equation}
\boxed{pV=RT.}
\end{equation}
This is the familiar ideal gas Law.
\end{node}

\subsection{Internal energy of ideal gas}

\begin{node}
We write internal energy as a function of volume and temperature
$e=e(V,T)$ to get the difference equation
\begin{equation}
  \D e = \left(\frac{\partial e}{\partial V}\right)_{T}\D V +
  \left(\frac{\partial e}{\partial T}\right)_{V}\D T.
\end{equation}
We use reciprocity Eq~\eqref{eq:second-reciprocity},
\begin{subequations}
\begin{equation}
\left(\frac{\partial e}{\partial V}\right)_{T} = T\left(\frac{\partial S}{\partial V}\right)_{T}-p,
\end{equation}
and a Maxwell relation for Helmholtz free energy Eq~\eqref{eq:maxwell-relation-for-helmholtz-free-energy},
\begin{equation}
\left(\frac{\partial S}{\partial V}\right)_{T} = \left(\frac{\partial p}{\partial T}\right)_{V},
\end{equation}
and the definition~(\S\ref{def:specific-heat-at-constant-volume}) of specific heat capacity at constant volume
\begin{equation}
c_{V} = \left(\frac{\partial e}{\partial T}\right)_{V} = T\left(\frac{\partial S}{\partial V}\right)_{T},
\end{equation}
\end{subequations}
to rewrite the difference equation as:
\begin{equation}\label{eq:fall:lec17:internal-energy:eq1}
\D e = \left[T\left(\frac{\partial p}{\partial T}\right)_{V}-p\right]\D p
+ c_{V}\,\D T.
\end{equation}
Then using the ideal gas Law $p=(R/V)T$ to find
\begin{equation}
\left(\frac{\partial p}{\partial T}\right)_{V}=\frac{R}{V}=\frac{p}{T},
\end{equation}
which kills off the bracketed term in Eq~\eqref{eq:fall:lec17:internal-energy:eq1},
giving us:
\begin{equation}\label{eq:fall:lec17:internal-energy:eq2}
\boxed{\D e = c_{V}\,\D T\quad\mbox{for ideal gases}.}
\end{equation}
This tells us that internal energy is a function of temperature
$e=e(T)$ for an ideal gas, and
\begin{equation}
\Delta e = e - e_{0} = \int^{T}_{T_{0}}c_{V}(T')\,\D T'\quad\mbox{for ideal gases},
\end{equation}
where $e_{0}$ is some reference internal energy (equivalently, $T_{0}$
is some reference temperature).
\end{node}

\subsection{Perfect Gases}

\begin{caution}
There are a variety of different ``perfect gases'' in the literature
(calorifically perfect gas, thermally perfect gas, etc.). Generically,
a ``perfect gas'' is an ideal gas plus some additional assumptions or
constraints.  
\end{caution}

\begin{node}
In a perfect gas, we assume $c_{p}$ and $c_{V}$ are constant. In this
case, we have
\begin{equation}
\Delta e = c_{V}\cdot(T-T_{0})=c_{V}\,\Delta T.
\end{equation}
\end{node}

\begin{node}
Perfect gases have no internal state, like Noble gases.
\end{node}

\begin{node}
In a non-perfect ideal gas, we can decompose the internal energy as
something like:
\begin{equation}
e = e_{\text{trans}} + e_{\text{rotational}} + e_{\text{vibration}}+e_{\text{electric}}+\dots
\end{equation}
which gives rise to a decomposition of specific heat at constant
volume
\begin{equation}
\left(\frac{\partial e}{\partial T}\right)_{V} = c_{V} = c_{V,\text{trans}}
+ c_{V,\text{rotational}} + c_{V,\text{vibration}} + c_{V,\text{electric}}+\dots.
\end{equation}
We can sketch what this might look like:
\begin{center}
\includegraphics{img/img.33}
\end{center}
For monatomic gases, the plot of $c_{V}/R$ may be drawn against
$T$. There is some threshold temperature $\Theta_{\text{Rot}}$ below
which rotation does not contribute to internal energy. Similarly,
there are thresholds $\Theta_{\text{vib}}$ for vibration and
$\Theta_{\text{dissoc}}$ for dissociation. Real gases are more
complicated. But this gives us an intuition for ideal non-perfect
gases. We can use step functions (as drawn), or smooth it out. One
common smooth variant
\begin{equation}
\frac{c_{V}}{R}=\frac{5}{2}+\frac{x^{2}\E^{x}}{(\E^{x}-1)^{2}}\quad\mbox{for $\Theta_{\text{Rot}}<T<\Theta_{\text{dissoc}}$},
\end{equation}
where $x=\Theta_{\text{vib}}/T$. [Note to self: Hydrogen is an example
  of this phenomenon, the graph plotted above.]
\end{node}

\begin{node}
Enthalpy for an ideal gas
\begin{subequations}
  \begin{align}
h &= e + \frac{p}{\rho}\quad\mbox{by definition of enthalpy}\\
&= e + pV\\
&= e(T) + RT\quad\mbox{by ideal gas Law},
  \end{align}
\end{subequations}
which means that enthalpy is a function of temperature alone $h=h(T)$
for ideal gases.
\end{node}

\begin{node}
The specific heat at constant pressure for ideal gases,
\begin{subequations}
  \begin{align}
c_{P} &= \left(\frac{\partial h}{\partial T}\right)_{p}\quad\mbox{by
  definition of $c_{p}$}\\
&= \left(\frac{\partial e}{\partial T}\right)_{p} + R\quad\mbox{for
  ideal gases}\\
&= c_{V} + R.
  \end{align}
\end{subequations}
So for ideal gases,
\begin{equation}
\boxed{c_{p}-c_{V}=R\quad\mbox{for ideal gases}.}
\end{equation}
\end{node}

\begin{node}
We similarly find the ratio of heat capacities $\gamma=c_{p}/c_{V}$
gives us
\begin{subequations}
  \begin{align}
    \gamma &= \frac{c_{p}}{c_{V}}\\
    &= \frac{c_{V}+R}{c_{V}}\quad\mbox{for ideal gases}\\
    &= 1 + \frac{R}{c_{V}}
  \end{align}
\end{subequations}
which can be rearranged to give us
\begin{equation}
\boxed{c_{V} = \frac{R}{\gamma-1}\quad\mbox{for ideal gases}.}
\end{equation}
Then we see after some manipulation that
\begin{equation}
\boxed{c_{p}=\frac{\gamma R}{\gamma-1}\quad\mbox{for ideal gases}.}
\end{equation}
\end{node}

\begin{node}
The speed of sound
\begin{subequations}
  \begin{align}
c^{2} &= \left(\frac{\partial p}{\partial\rho}\right)_{S}\quad\mbox{by definition}\\
&= -V^{2}\left(\frac{\partial p}{\partial V}\right)_{S}\quad\mbox{for ideal gases}\\
&= - V^{2}\gamma\left(\frac{\partial p}{\partial V}\right)\quad\mbox{using $\kappa_{T}=\gamma\kappa_{S}$}\\
&= -V^{2}\gamma\left(\frac{-RT}{V^{2}}\right)\quad\mbox{using $p=RT/V$}.
  \end{align}
\end{subequations}
Hence
\begin{equation}
\boxed{c^{2}=\gamma RT = \gamma\frac{p}{\rho}\quad\mbox{for ideal gases}.}
\end{equation}
\end{node}

\begin{remark}
For ideal gases, we find many definitions which required two
parameters simplify considerably and now require only one parameter.
\end{remark}
